% ===========================================================================
\section{Arithmetizing the ECDSA building blocks}\label{sec:arith}
% ===========================================================================

We collect the arithmetic identities that translate each step of the
double-and-add loop into a constraint over $\F_p$. Throughout this
section, all arithmetic is in $\F_p$ unless noted otherwise; we
suppress reductions $\bmod\,p$.

\subsection{Shamir's trick}\label{sec:shamir}

The standard binary double-and-add algorithm computes a single
scalar multiplication $u A$ (for a $k$-bit scalar $u$ and a point $A$)
by initializing an accumulator $P \leftarrow \OOO$ and, for
$j = k-1, \dots, 0$, performing
$$
P \;\leftarrow\; 2P, \qquad
P \;\leftarrow\; P + A \quad \text{whenever bit $j$ of $u$ is $1$.}
$$
Naively, computing the two-term sum $u_1 G + u_2 Q$ would run two
such loops for a total of $2k$ doublings and up to $2k$ additions.
Shamir's trick merges them into a single loop. Precompute the four
points
$$
T_{00} = \OOO,\quad T_{10} = G,\quad T_{01} = Q,\quad T_{11} = G + Q,
$$
initialize $P \leftarrow \OOO$, and run, for $j = k-1, \dots, 0$,
$$
P \;\leftarrow\; 2P, \qquad
P \;\leftarrow\; P + T_{b_1[j]\,b_2[j]}.
$$
After $k = 256$ steps, $P = u_1 G + u_2 Q$. The total cost is $k$
doublings and at most $k$ mixed additions.

In the UAIR, the per-row affine addend
$T_{b_1[t]\,b_2[t]} \in \{\OOO, G, Q, G+Q\}$ is selected
\emph{in-circuit} from the bit pair via the linear formula
\begin{equation}\label{eq:addend-formula}
  T \;=\; b_1\,(1 - b_2)\,G \;+\; (1 - b_1)\,b_2\,Q \;+\; b_1 b_2\,(G + Q),
\end{equation}
applied coefficient-wise. Equivalently, expanded into the canonical
``affine + $b_1 b_2$ correction'' form,
\begin{equation}\label{eq:addend-formula-expanded}
  T \;=\; b_1 G \;+\; b_2 Q \;+\; b_1 b_2 \bigl((G+Q) - G - Q\bigr).
\end{equation}
Five public columns carry the inputs to \eqref{eq:addend-formula-expanded}:
$(\col{PA\_B1}, \col{PA\_B2})$ for the bit pair $(b_1, b_2)$;
$(\col{PA\_QX}, \col{PA\_QY})$ for the affine coordinates of the
public-key point $Q$; and $(\col{PA\_QGX}, \col{PA\_QGY})$ for the
verifier-derivable affine coordinates of $G + Q$. The generator $G$
enters as a UAIR scalar (constant across all proofs of this UAIR).
The boolean $\col{S\_ADD}$ remains a public column with the relation
$\col{S\_ADD}[t] = b_1[t] + b_2[t] - b_1[t] b_2[t]$ (i.e.\ the
non-identity flag); when $\col{S\_ADD}[t] = 0$ the addition step at
row $t$ collapses to $P \leftarrow 2P$.

The expression \eqref{eq:addend-formula-expanded} is \emph{inlined}
directly inside the addition-scratch constraints \eqref{eq:H} and
\eqref{eq:Ra} of \Cref{sec:add} below — there is no separate
$T_X$ or $T_Y$ witness column, and no separate ``addend
selection'' constraint family.

\begin{remark}[Why in-circuit, and what it costs]\label{rem:addend-cost}
A natural alternative is to have the verifier write the chosen
affine addend directly into $(\col{PA\_X\_ADDEND}, \col{PA\_Y\_ADDEND})$
for each row, exactly as in a precomputed $256$-row table lookup
performed off-protocol from the bits. The deployed design
deliberately rejects that route to keep the proof self-contained —
i.e., to avoid a hidden ``verifier oracle'' that materialises a
non-trivial computation step (the $4$-way table lookup driven by
public bits) outside the in-circuit checks.

The cost of the chosen in-circuit design relative to the
verifier-supplied alternative:
\begin{itemize}[leftmargin=2em]
  \item \textbf{$+4$ public columns.} The deployed layout drops
    $\col{PA\_X\_ADDEND}$ and $\col{PA\_Y\_ADDEND}$ ($-2$) and adds
    $\col{PA\_B1}, \col{PA\_B2}, \col{PA\_QX}, \col{PA\_QY},
    \col{PA\_QGX}, \col{PA\_QGY}$ ($+6$). The two $Q$- and
    $(G+Q)$-coordinate pairs are constant across rows of a single
    proof.
  \item \textbf{Constraint-degree elevation, $6 \to 7$.} The expression
    \eqref{eq:addend-formula-expanded} is degree~$3$ in trace cells
    (the $b_1 b_2 \cdot (G+Q)$ term reaches three columns). Inlined
    into $H = T_x P_{\mathrm{pa},Z}^2 - P_{\mathrm{pa},X}$ multiplied
    against $\col{S\_ACTIVE}$, the C-A1 constraint reaches degree~$6$
    (was~$4$); the C-A2 constraint with $Z_{\mathrm{pa}}^3$ reaches
    degree~$7$ (was~$5$). The previous design's max constraint
    degree was~$6$ from O2 (the Y output-selection); the deployed
    UAIR's max is~$7$, set by C-A2.
  \item \textbf{$0$ extra in-circuit constraints.} The selection
    formula is inlined; no separate ``$T_X = \cdots$'' equation is
    posted, and no booleanity constraints are emitted.
\end{itemize}

\paragraph{Soundness obligation (verifier-side, not yet enforced).}
The $\F_p$-typed public columns $\col{PA\_B1}, \col{PA\_B2}$ are
\emph{not} natively constrained to $\{0,1\}$ by the framework
(integer columns admit any $320$-bit canonical representative). A
malicious prover could place arbitrary values in them and compose any
addend it likes via \eqref{eq:addend-formula-expanded}. The intended
discharge is the verifier-side \texttt{verify\_public\_structure}
trait method, following the same row-wise inspection pattern the
SHA UAIR uses for compensator-zero checks; concretely, the verifier
must check
$\col{PA\_B1}[t], \col{PA\_B2}[t] \in \{0, 1\}$ at every active row,
$\col{S\_ADD}[t] = b_1[t] + b_2[t] - b_1[t] b_2[t]$, and
$(\col{PA\_QX}, \col{PA\_QY}, \col{PA\_QGX}, \col{PA\_QGY})$ constant
across rows of a single proof. The first two checks are not
expressible under the current
\texttt{Uair::verify\_public\_structure} bound
($\mathtt{IntT}: \mathtt{Clone} + \mathtt{num\_traits::Zero}$) — they
need a $\mathtt{PartialEq} + \mathtt{One}$ widening of the bound.
The implementation defers this to a follow-up; the synthetic
end-to-end test populates the bit columns honestly.
\end{remark}

\subsection{Jacobian doubling}\label{sec:doubling}

Given an input Jacobian point $\jac{P}$, the Jacobian doubling
$2P = (P_{\mathrm{pa},X} : P_{\mathrm{pa},Y} : P_{\mathrm{pa},Z})$ is
computed in $\F_p$ by
\begin{align}
P_{\mathrm{pa},Z} &= 2 P_Y P_Z, \label{eq:dbl-z}\\
P_{\mathrm{pa},X} &= 9 P_X^4 - 8 P_X P_Y^2, \label{eq:dbl-x}\\
P_{\mathrm{pa},Y} &= 12 P_X^3 P_Y^2 - 3 P_X^2 P_{\mathrm{pa},X} - 8 P_Y^4. \label{eq:dbl-y}
\end{align}
These are the standard short-Weierstrass $a = 0$ doubling formulas.
The intermediate $S = P_Y^2$ is computed inline (the implementation
inlines it rather than committing it as a separate column).

\subsection{Mixed Jacobian-affine addition}\label{sec:add}

Given $(P_{\mathrm{pa},X} : P_{\mathrm{pa},Y} : P_{\mathrm{pa},Z})$
(the doubled point) and an affine addend $(T_X, T_Y)$ from the
precomputed $\{G, Q, G+Q\}$, the mixed addition
$P_{\mathrm{pa}} + T = (P_{\mathrm{add},X} : P_{\mathrm{add},Y} : P_{\mathrm{add},Z})$
proceeds via the scratch
quantities
\begin{align}
H   &:= T_X \cdot P_{\mathrm{pa},Z}^2 \;-\; P_{\mathrm{pa},X}, \label{eq:H}\\
R_a &:= T_Y \cdot P_{\mathrm{pa},Z}^3 \;-\; P_{\mathrm{pa},Y}, \label{eq:Ra}
\end{align}
followed by the addition outputs
\begin{align}
P_{\mathrm{add},Z} &= P_{\mathrm{pa},Z} \cdot H, \label{eq:add-z}\\
P_{\mathrm{add},X} &= R_a^2 - H^3 - 2 P_{\mathrm{pa},X} H^2, \label{eq:add-x}\\
P_{\mathrm{add},Y} &= R_a \bigl(P_{\mathrm{pa},X} H^2 - P_{\mathrm{add},X}\bigr) - P_{\mathrm{pa},Y} H^3. \label{eq:add-y}
\end{align}
The implementation commits only $H$ (named $\col{W\_C}$) and $R_a$
(named $\col{W\_D}$), inlining $P_{\mathrm{pa},Z}^2$, $P_{\mathrm{pa},Z}^3$,
$H^2$, $H^3$, and $P_{\mathrm{pa},X} H^2$ into the surrounding
constraint expressions. This trades higher per-constraint algebraic
degree for fewer committed columns; with the inlined form the maximum
constraint degree is~$6$ (attained in \eqref{eq:add-y} after
selector gating, see \Cref{sec:constraints}).

\subsection{Conditional addition via output-selection}\label{sec:cond-add}

Rather than gating the addition step with a multiplicative selector
that would inflate the constraint degree, the implementation folds
the conditional add into the row-chaining itself. Let
$\col{S\_ADD}[t] \in \{0, 1\}$ be the per-row addition selector
(\Cref{rem:addend-cost}). The next-row Jacobian state is then defined as
\begin{align}
P_{\mathrm{next},X} &= P_{\mathrm{pa},X} + \col{S\_ADD}\bigl(P_{\mathrm{add},X} - P_{\mathrm{pa},X}\bigr), \label{eq:next-x}\\
P_{\mathrm{next},Y} &= P_{\mathrm{pa},Y} + \col{S\_ADD}\bigl(P_{\mathrm{add},Y} - P_{\mathrm{pa},Y}\bigr), \label{eq:next-y}\\
P_{\mathrm{next},Z} &= P_{\mathrm{pa},Z} + \col{S\_ADD}\bigl(P_{\mathrm{add},Z} - P_{\mathrm{pa},Z}\bigr), \label{eq:next-z}
\end{align}
so that $\col{S\_ADD} = 1$ selects $P_{\mathrm{add}}$ and
$\col{S\_ADD} = 0$ selects $P_{\mathrm{pa}}$ (the pure doubling result).

Because $\col{S\_ADD}$ is a public boolean, $\col{S\_ADD} (1 - \col{S\_ADD}) = 0$
need not be enforced in-circuit; the verifier supplies the column with
$0/1$ values directly.

\begin{remark}[Inlined addition collapses to a single output-selection constraint]
\label{rem:inlined-add}
Substituting $P_{\mathrm{add}, X/Y/Z}$ from
\eqref{eq:add-z}--\eqref{eq:add-y} into \eqref{eq:next-x}--\eqref{eq:next-z}
and inlining all of $P_{\mathrm{pa},Z}^2, H^2, H^3, P_{\mathrm{pa},X} H^2$
yields three constraints whose maximum algebraic degree (in committed
columns) is~$6$. The deg-$6$ monomial in the $Y$ output-selection
constraint is exactly $\col{S\_ACTIVE} \cdot \col{S\_ADD} \cdot D \cdot X_{\mathrm{pa}} \cdot C^2$,
matching the spec's $s_{\mathrm{reg}} \cdot R_a \cdot X_{\mathrm{mid}} \cdot H^2$
term.
\end{remark}

\subsection{Initial step and final state}\label{sec:boundary}

\paragraph{Init boundary.}
At row~$0$ the Jacobian state is pinned to the verifier-supplied
non-identity seed $R_{\mathrm{init}} = (R_{\mathrm{init},X} : R_{\mathrm{init},Y} : R_{\mathrm{init},Z})$
(\Cref{rem:r-init}):
$$
\bigl(P_X[0], P_Y[0], P_Z[0]\bigr) \;=\; \bigl(\col{PA\_R\_INIT\_X}, \col{PA\_R\_INIT\_Y}, \col{PA\_R\_INIT\_Z}\bigr).
$$
This is enforced by three zero-ideal constraints gated by
$\col{S\_INIT}$ (see \Cref{con:init-x}--\Cref{con:init-z}).

\paragraph{Final-row affine readout (in-circuit).}
At the final row $t = \mathrm{FINAL\_ROW} := \mathrm{NUM\_SHAMIR\_ROUNDS} = 256$,
the witness Jacobian state $\jac{P}[256]$ holds the result
$R = u_1 G + u_2 Q$ (modulo the $R_{\mathrm{init}}$ offset of
\Cref{rem:r-init}). The deployed UAIR extracts the affine
$x$-coordinate of $R$ in-circuit using two public columns and two
$\col{S\_FINAL}$-gated $(p)$-ideal constraints:
\begin{align}
\col{S\_FINAL}[t] \cdot \bigl(P_Z[t] \cdot \col{PA\_Z\_INV}[t] - 1\bigr) &\in (p), \tag{F1}\label{con:f1}\\
\col{S\_FINAL}[t] \cdot \bigl(P_X[t] \cdot \col{PA\_Z\_INV}[t]^2 - \col{PA\_R\_X}[t]\bigr) &\in (p). \tag{F2}\label{con:f2}
\end{align}
Constraint \ref{con:f1} forces $P_Z[256] \neq 0$ (otherwise the equation
is unsatisfiable, ruling out the point at infinity) and pins
$\col{PA\_Z\_INV}[256]$ to the multiplicative inverse of $P_Z[256]$
modulo $p$. Constraint \ref{con:f2} then pins $\col{PA\_R\_X}[256]$
to the affine $x$-coordinate $P_X[256] \cdot \col{PA\_Z\_INV}[256]^2$.
Both columns are public; their cells at rows $t \neq 256$ are
unconstrained (the gating selector zeroes the constraint there) and
are populated with $0$ by the trace builder.

\begin{remark}[The order-field check is left off-protocol]\label{rem:final-readout}
The signature-defining identity $R_x \equiv r \pmod n$ — where $n$ is
the curve order, distinct from the base prime $p$ — is \emph{not}
enforced in-circuit. Because both $\col{PA\_R\_X}$ and $r$ are
public ($\col{PA\_R\_X}$ is part of \texttt{public\_trace}, $r$ is
part of the verifier's input alongside the message and signature),
the verifier checks the equality directly off-protocol as a single
$\F_n$-arithmetic comparison after \texttt{verify} returns; no
in-circuit constraint or extra column is required for this step.
The alternative would be to add $n$ to the index, turning
$\mathbf{q} = (p)$ into $\mathbf{q} = (p, n)$, with an $(n)$-ideal
constraint $\col{S\_FINAL} \cdot (\col{PA\_R\_X}[t] - r[t]) \in (n)$.
The current single-prime design avoids the multi-prime branch and
relies on the fact that the order-field check is a public-public
equality, which the protocol layer above the UAIR can discharge
directly.
\end{remark}
